% Options for packages loaded elsewhere
\PassOptionsToPackage{unicode}{hyperref}
\PassOptionsToPackage{hyphens}{url}
\PassOptionsToPackage{dvipsnames,svgnames,x11names}{xcolor}
%
\documentclass[
  letterpaper,
  DIV=11,
  numbers=noendperiod,
  oneside]{scrartcl}

\usepackage{amsmath,amssymb}
\usepackage{iftex}
\ifPDFTeX
  \usepackage[T1]{fontenc}
  \usepackage[utf8]{inputenc}
  \usepackage{textcomp} % provide euro and other symbols
\else % if luatex or xetex
  \usepackage{unicode-math}
  \defaultfontfeatures{Scale=MatchLowercase}
  \defaultfontfeatures[\rmfamily]{Ligatures=TeX,Scale=1}
\fi
\usepackage{lmodern}
\ifPDFTeX\else  
    % xetex/luatex font selection
\fi
% Use upquote if available, for straight quotes in verbatim environments
\IfFileExists{upquote.sty}{\usepackage{upquote}}{}
\IfFileExists{microtype.sty}{% use microtype if available
  \usepackage[]{microtype}
  \UseMicrotypeSet[protrusion]{basicmath} % disable protrusion for tt fonts
}{}
\makeatletter
\@ifundefined{KOMAClassName}{% if non-KOMA class
  \IfFileExists{parskip.sty}{%
    \usepackage{parskip}
  }{% else
    \setlength{\parindent}{0pt}
    \setlength{\parskip}{6pt plus 2pt minus 1pt}}
}{% if KOMA class
  \KOMAoptions{parskip=half}}
\makeatother
\usepackage{xcolor}
\usepackage[left=1in,marginparwidth=2.0666666666667in,textwidth=4.1333333333333in,marginparsep=0.3in]{geometry}
\setlength{\emergencystretch}{3em} % prevent overfull lines
\setcounter{secnumdepth}{-\maxdimen} % remove section numbering
% Make \paragraph and \subparagraph free-standing
\makeatletter
\ifx\paragraph\undefined\else
  \let\oldparagraph\paragraph
  \renewcommand{\paragraph}{
    \@ifstar
      \xxxParagraphStar
      \xxxParagraphNoStar
  }
  \newcommand{\xxxParagraphStar}[1]{\oldparagraph*{#1}\mbox{}}
  \newcommand{\xxxParagraphNoStar}[1]{\oldparagraph{#1}\mbox{}}
\fi
\ifx\subparagraph\undefined\else
  \let\oldsubparagraph\subparagraph
  \renewcommand{\subparagraph}{
    \@ifstar
      \xxxSubParagraphStar
      \xxxSubParagraphNoStar
  }
  \newcommand{\xxxSubParagraphStar}[1]{\oldsubparagraph*{#1}\mbox{}}
  \newcommand{\xxxSubParagraphNoStar}[1]{\oldsubparagraph{#1}\mbox{}}
\fi
\makeatother

\usepackage{color}
\usepackage{fancyvrb}
\newcommand{\VerbBar}{|}
\newcommand{\VERB}{\Verb[commandchars=\\\{\}]}
\DefineVerbatimEnvironment{Highlighting}{Verbatim}{commandchars=\\\{\}}
% Add ',fontsize=\small' for more characters per line
\usepackage{framed}
\definecolor{shadecolor}{RGB}{241,243,245}
\newenvironment{Shaded}{\begin{snugshade}}{\end{snugshade}}
\newcommand{\AlertTok}[1]{\textcolor[rgb]{0.68,0.00,0.00}{#1}}
\newcommand{\AnnotationTok}[1]{\textcolor[rgb]{0.37,0.37,0.37}{#1}}
\newcommand{\AttributeTok}[1]{\textcolor[rgb]{0.40,0.45,0.13}{#1}}
\newcommand{\BaseNTok}[1]{\textcolor[rgb]{0.68,0.00,0.00}{#1}}
\newcommand{\BuiltInTok}[1]{\textcolor[rgb]{0.00,0.23,0.31}{#1}}
\newcommand{\CharTok}[1]{\textcolor[rgb]{0.13,0.47,0.30}{#1}}
\newcommand{\CommentTok}[1]{\textcolor[rgb]{0.37,0.37,0.37}{#1}}
\newcommand{\CommentVarTok}[1]{\textcolor[rgb]{0.37,0.37,0.37}{\textit{#1}}}
\newcommand{\ConstantTok}[1]{\textcolor[rgb]{0.56,0.35,0.01}{#1}}
\newcommand{\ControlFlowTok}[1]{\textcolor[rgb]{0.00,0.23,0.31}{\textbf{#1}}}
\newcommand{\DataTypeTok}[1]{\textcolor[rgb]{0.68,0.00,0.00}{#1}}
\newcommand{\DecValTok}[1]{\textcolor[rgb]{0.68,0.00,0.00}{#1}}
\newcommand{\DocumentationTok}[1]{\textcolor[rgb]{0.37,0.37,0.37}{\textit{#1}}}
\newcommand{\ErrorTok}[1]{\textcolor[rgb]{0.68,0.00,0.00}{#1}}
\newcommand{\ExtensionTok}[1]{\textcolor[rgb]{0.00,0.23,0.31}{#1}}
\newcommand{\FloatTok}[1]{\textcolor[rgb]{0.68,0.00,0.00}{#1}}
\newcommand{\FunctionTok}[1]{\textcolor[rgb]{0.28,0.35,0.67}{#1}}
\newcommand{\ImportTok}[1]{\textcolor[rgb]{0.00,0.46,0.62}{#1}}
\newcommand{\InformationTok}[1]{\textcolor[rgb]{0.37,0.37,0.37}{#1}}
\newcommand{\KeywordTok}[1]{\textcolor[rgb]{0.00,0.23,0.31}{\textbf{#1}}}
\newcommand{\NormalTok}[1]{\textcolor[rgb]{0.00,0.23,0.31}{#1}}
\newcommand{\OperatorTok}[1]{\textcolor[rgb]{0.37,0.37,0.37}{#1}}
\newcommand{\OtherTok}[1]{\textcolor[rgb]{0.00,0.23,0.31}{#1}}
\newcommand{\PreprocessorTok}[1]{\textcolor[rgb]{0.68,0.00,0.00}{#1}}
\newcommand{\RegionMarkerTok}[1]{\textcolor[rgb]{0.00,0.23,0.31}{#1}}
\newcommand{\SpecialCharTok}[1]{\textcolor[rgb]{0.37,0.37,0.37}{#1}}
\newcommand{\SpecialStringTok}[1]{\textcolor[rgb]{0.13,0.47,0.30}{#1}}
\newcommand{\StringTok}[1]{\textcolor[rgb]{0.13,0.47,0.30}{#1}}
\newcommand{\VariableTok}[1]{\textcolor[rgb]{0.07,0.07,0.07}{#1}}
\newcommand{\VerbatimStringTok}[1]{\textcolor[rgb]{0.13,0.47,0.30}{#1}}
\newcommand{\WarningTok}[1]{\textcolor[rgb]{0.37,0.37,0.37}{\textit{#1}}}

\providecommand{\tightlist}{%
  \setlength{\itemsep}{0pt}\setlength{\parskip}{0pt}}\usepackage{longtable,booktabs,array}
\usepackage{calc} % for calculating minipage widths
% Correct order of tables after \paragraph or \subparagraph
\usepackage{etoolbox}
\makeatletter
\patchcmd\longtable{\par}{\if@noskipsec\mbox{}\fi\par}{}{}
\makeatother
% Allow footnotes in longtable head/foot
\IfFileExists{footnotehyper.sty}{\usepackage{footnotehyper}}{\usepackage{footnote}}
\makesavenoteenv{longtable}
\usepackage{graphicx}
\makeatletter
\def\maxwidth{\ifdim\Gin@nat@width>\linewidth\linewidth\else\Gin@nat@width\fi}
\def\maxheight{\ifdim\Gin@nat@height>\textheight\textheight\else\Gin@nat@height\fi}
\makeatother
% Scale images if necessary, so that they will not overflow the page
% margins by default, and it is still possible to overwrite the defaults
% using explicit options in \includegraphics[width, height, ...]{}
\setkeys{Gin}{width=\maxwidth,height=\maxheight,keepaspectratio}
% Set default figure placement to htbp
\makeatletter
\def\fps@figure{htbp}
\makeatother

\KOMAoption{captions}{tableheading}
\makeatletter
\@ifpackageloaded{caption}{}{\usepackage{caption}}
\AtBeginDocument{%
\ifdefined\contentsname
  \renewcommand*\contentsname{Table of contents}
\else
  \newcommand\contentsname{Table of contents}
\fi
\ifdefined\listfigurename
  \renewcommand*\listfigurename{List of Figures}
\else
  \newcommand\listfigurename{List of Figures}
\fi
\ifdefined\listtablename
  \renewcommand*\listtablename{List of Tables}
\else
  \newcommand\listtablename{List of Tables}
\fi
\ifdefined\figurename
  \renewcommand*\figurename{Figure}
\else
  \newcommand\figurename{Figure}
\fi
\ifdefined\tablename
  \renewcommand*\tablename{Table}
\else
  \newcommand\tablename{Table}
\fi
}
\@ifpackageloaded{float}{}{\usepackage{float}}
\floatstyle{ruled}
\@ifundefined{c@chapter}{\newfloat{codelisting}{h}{lop}}{\newfloat{codelisting}{h}{lop}[chapter]}
\floatname{codelisting}{Listing}
\newcommand*\listoflistings{\listof{codelisting}{List of Listings}}
\makeatother
\makeatletter
\makeatother
\makeatletter
\@ifpackageloaded{caption}{}{\usepackage{caption}}
\@ifpackageloaded{subcaption}{}{\usepackage{subcaption}}
\makeatother
\makeatletter
\@ifpackageloaded{sidenotes}{}{\usepackage{sidenotes}}
\@ifpackageloaded{marginnote}{}{\usepackage{marginnote}}
\makeatother

\ifLuaTeX
  \usepackage{selnolig}  % disable illegal ligatures
\fi
\usepackage{bookmark}

\IfFileExists{xurl.sty}{\usepackage{xurl}}{} % add URL line breaks if available
\urlstyle{same} % disable monospaced font for URLs
\hypersetup{
  pdftitle={Prediction with (shallow) Neural Networks},
  colorlinks=true,
  linkcolor={blue},
  filecolor={Maroon},
  citecolor={Blue},
  urlcolor={Blue},
  pdfcreator={LaTeX via pandoc}}


\title{Prediction with (shallow) Neural Networks}
\author{Adapted by EVL, FRC and ASP}
\date{2025-03-17}

\begin{document}
\maketitle


\section{An example using R}\label{an-example-using-r}

\subsection{A predictive ANN}\label{a-predictive-ann}

We use the \texttt{neuralnet} package to build a simple neural network
to predict if a type of stock pays dividends or not.

\begin{Shaded}
\begin{Highlighting}[]
\ControlFlowTok{if}\NormalTok{ (}\SpecialCharTok{!}\FunctionTok{require}\NormalTok{(neuralnet)) }
  \FunctionTok{install.packages}\NormalTok{(}\StringTok{"neuralnet"}\NormalTok{, }\AttributeTok{dep=}\ConstantTok{TRUE}\NormalTok{)}
\end{Highlighting}
\end{Shaded}

\subsection{Data for the example}\label{data-for-the-example}

And use the \texttt{dividendinfo.csv} dataset from
\url{https://github.com/MGCodesandStats/datasets}

\begin{Shaded}
\begin{Highlighting}[]
\NormalTok{mydata }\OtherTok{\textless{}{-}} \FunctionTok{read.csv}\NormalTok{(}\StringTok{"https://raw.githubusercontent.com/MGCodesandStats/datasets/master/dividendinfo.csv"}\NormalTok{)}
\FunctionTok{str}\NormalTok{(mydata)}
\end{Highlighting}
\end{Shaded}

\begin{verbatim}
'data.frame':   200 obs. of  6 variables:
 $ dividend       : int  0 1 1 0 1 1 1 0 1 1 ...
 $ fcfps          : num  2.75 4.96 2.78 0.43 2.94 3.9 1.09 2.32 2.5 4.46 ...
 $ earnings_growth: num  -19.25 0.83 1.09 12.97 2.44 ...
 $ de             : num  1.11 1.09 0.19 1.7 1.83 0.46 2.32 3.34 3.15 3.33 ...
 $ mcap           : int  545 630 562 388 684 621 656 351 658 330 ...
 $ current_ratio  : num  0.924 1.469 1.976 1.942 2.487 ...
\end{verbatim}

\subsection{Data pre-processing}\label{data-pre-processing}

\begin{Shaded}
\begin{Highlighting}[]
\NormalTok{normalize }\OtherTok{\textless{}{-}} \ControlFlowTok{function}\NormalTok{(x) \{}
  \FunctionTok{return}\NormalTok{ ((x }\SpecialCharTok{{-}} \FunctionTok{min}\NormalTok{(x)) }\SpecialCharTok{/}\NormalTok{ (}\FunctionTok{max}\NormalTok{(x) }\SpecialCharTok{{-}} \FunctionTok{min}\NormalTok{(x)))}
\NormalTok{\}}
\NormalTok{normData }\OtherTok{\textless{}{-}} \FunctionTok{as.data.frame}\NormalTok{(}\FunctionTok{lapply}\NormalTok{(mydata, normalize))}
\end{Highlighting}
\end{Shaded}

\subsection{Test and training sets}\label{test-and-training-sets}

Finally we break our data in a test and a training set:

\begin{Shaded}
\begin{Highlighting}[]
\NormalTok{perc2Train }\OtherTok{\textless{}{-}} \DecValTok{2}\SpecialCharTok{/}\DecValTok{3}
\NormalTok{ssize }\OtherTok{\textless{}{-}} \FunctionTok{nrow}\NormalTok{(normData)}
\FunctionTok{set.seed}\NormalTok{(}\DecValTok{12345}\NormalTok{)}
\NormalTok{data\_rows }\OtherTok{\textless{}{-}} \FunctionTok{floor}\NormalTok{(perc2Train }\SpecialCharTok{*}\NormalTok{ssize)}
\NormalTok{train\_indices }\OtherTok{\textless{}{-}} \FunctionTok{sample}\NormalTok{(}\FunctionTok{c}\NormalTok{(}\DecValTok{1}\SpecialCharTok{:}\NormalTok{ssize), data\_rows)}
\NormalTok{trainset }\OtherTok{\textless{}{-}}\NormalTok{ normData[train\_indices,]}
\NormalTok{testset }\OtherTok{\textless{}{-}}\NormalTok{ normData[}\SpecialCharTok{{-}}\NormalTok{train\_indices,]}
\end{Highlighting}
\end{Shaded}

\subsection{Training a neural network}\label{training-a-neural-network}

We train a simple NN with two hidden layers, with 4 and 2 neurons
respectively.

First we deine the parameters and hyperparameters

\begin{Shaded}
\begin{Highlighting}[]
\NormalTok{hidden\_layers }\OtherTok{\textless{}{-}} \FunctionTok{c}\NormalTok{(}\DecValTok{4}\NormalTok{,}\DecValTok{2}\NormalTok{)  }\CommentTok{\# Two hidden layers with 4 and 2 neurons}
\NormalTok{activation\_function }\OtherTok{\textless{}{-}} \StringTok{"logistic"}  \CommentTok{\# Sigmoid}
\NormalTok{optimizer }\OtherTok{\textless{}{-}} \StringTok{"rprop+"}  \CommentTok{\# Optimizer}
\NormalTok{learning\_rate }\OtherTok{\textless{}{-}} \FloatTok{0.01}  \CommentTok{\# Learning rate}
\NormalTok{threshold\_value }\OtherTok{\textless{}{-}} \FloatTok{0.01}  \CommentTok{\# Convergence  thresholod }
\NormalTok{max\_iterations }\OtherTok{\textless{}{-}} \FloatTok{1e6}  \CommentTok{\# Max iterations}
\end{Highlighting}
\end{Shaded}

\begin{Shaded}
\begin{Highlighting}[]
\CommentTok{\# Neural Network}
\FunctionTok{library}\NormalTok{(neuralnet)}
\NormalTok{nn }\OtherTok{\textless{}{-}} \FunctionTok{neuralnet}\NormalTok{(}\CommentTok{\#dividend \textasciitilde{} .,}
\NormalTok{                dividend }\SpecialCharTok{\textasciitilde{}}\NormalTok{ fcfps }\SpecialCharTok{+}\NormalTok{ earnings\_growth }\SpecialCharTok{+}\NormalTok{ de }\SpecialCharTok{+}\NormalTok{ mcap }\SpecialCharTok{+}\NormalTok{ current\_ratio, }
                \AttributeTok{data=}\NormalTok{trainset, }
                \AttributeTok{hidden=}\NormalTok{hidden\_layers, }
                \AttributeTok{act.fct=}\NormalTok{activation\_function, }
                \AttributeTok{algorithm=}\NormalTok{optimizer, }
                \AttributeTok{learningrate=}\NormalTok{learning\_rate,}
                \AttributeTok{linear.output=}\ConstantTok{FALSE}\NormalTok{, }
                \AttributeTok{threshold=}\NormalTok{threshold\_value,}
                \AttributeTok{stepmax=}\NormalTok{max\_iterations)}
\end{Highlighting}
\end{Shaded}

\subsection{Network plot}\label{network-plot}

The output of the procedure is a neural network with estimated weights.

\begin{Shaded}
\begin{Highlighting}[]
\FunctionTok{plot}\NormalTok{(nn, }\AttributeTok{rep =} \StringTok{"best"}\NormalTok{)}
\end{Highlighting}
\end{Shaded}

\includegraphics{Lab-C3.2-DividendPrediction_ANN_files/figure-pdf/unnamed-chunk-7-1.pdf}

In order to facilitate visualization we can round decimals to 2 for all
weights.

\begin{Shaded}
\begin{Highlighting}[]
\NormalTok{round\_weights }\OtherTok{\textless{}{-}} \ControlFlowTok{function}\NormalTok{(nn\_model, }\AttributeTok{decimals=}\DecValTok{2}\NormalTok{) \{}
\NormalTok{  nn\_model}\SpecialCharTok{$}\NormalTok{weights }\OtherTok{\textless{}{-}} \FunctionTok{lapply}\NormalTok{(nn\_model}\SpecialCharTok{$}\NormalTok{weights, }\ControlFlowTok{function}\NormalTok{(layer) \{}
    \FunctionTok{lapply}\NormalTok{(layer, }\ControlFlowTok{function}\NormalTok{(weight\_matrix) \{}
      \ControlFlowTok{if}\NormalTok{ (}\FunctionTok{is.numeric}\NormalTok{(weight\_matrix)) \{}
        \FunctionTok{round}\NormalTok{(weight\_matrix, decimals)}
\NormalTok{      \} }\ControlFlowTok{else}\NormalTok{ \{}
\NormalTok{        weight\_matrix}
\NormalTok{      \}}
\NormalTok{    \})}
\NormalTok{  \})}
  \FunctionTok{return}\NormalTok{(nn\_model)}
\NormalTok{\}}

\CommentTok{\# Apply function to trained net}
\NormalTok{nn\_rounded }\OtherTok{\textless{}{-}} \FunctionTok{round\_weights}\NormalTok{(nn, }\DecValTok{2}\NormalTok{)}

\CommentTok{\# View network with less decimals}
\FunctionTok{plot}\NormalTok{(nn\_rounded, }\AttributeTok{rep =} \StringTok{"best"}\NormalTok{)}
\end{Highlighting}
\end{Shaded}

\includegraphics{Lab-C3.2-DividendPrediction_ANN_files/figure-pdf/unnamed-chunk-8-1.pdf}

\subsection{Predictions}\label{predictions}

\begin{Shaded}
\begin{Highlighting}[]
\NormalTok{temp\_test }\OtherTok{\textless{}{-}} \FunctionTok{subset}\NormalTok{(testset, }\AttributeTok{select =}
                      \FunctionTok{c}\NormalTok{(}\StringTok{"fcfps"}\NormalTok{,}\StringTok{"earnings\_growth"}\NormalTok{, }
                        \StringTok{"de"}\NormalTok{, }\StringTok{"mcap"}\NormalTok{, }\StringTok{"current\_ratio"}\NormalTok{))}
\NormalTok{nn.results }\OtherTok{\textless{}{-}} \FunctionTok{compute}\NormalTok{(nn, temp\_test)}
\NormalTok{results }\OtherTok{\textless{}{-}} \FunctionTok{data.frame}\NormalTok{(}\AttributeTok{actual =} 
\NormalTok{                  testset}\SpecialCharTok{$}\NormalTok{dividend, }
                  \AttributeTok{prediction =}\NormalTok{ nn.results}\SpecialCharTok{$}\NormalTok{net.result)}
\FunctionTok{head}\NormalTok{(results)}
\end{Highlighting}
\end{Shaded}

\begin{verbatim}
   actual   prediction
9       1 1.000000e+00
19      1 9.999167e-01
22      0 1.184302e-26
26      0 4.282788e-09
27      1 3.686182e-03
29      1 9.999999e-01
\end{verbatim}

\subsection{Model evaluation}\label{model-evaluation}

\begin{Shaded}
\begin{Highlighting}[]
\NormalTok{roundedresults}\OtherTok{\textless{}{-}}\FunctionTok{sapply}\NormalTok{(results,round,}\AttributeTok{digits=}\DecValTok{0}\NormalTok{)}
\NormalTok{roundedresultsdf}\OtherTok{=}\FunctionTok{data.frame}\NormalTok{(roundedresults)}
\FunctionTok{attach}\NormalTok{(roundedresultsdf)}
\FunctionTok{table}\NormalTok{(actual,prediction)}
\end{Highlighting}
\end{Shaded}

\begin{verbatim}
      prediction
actual  0  1
     0 35  1
     1  7 24
\end{verbatim}

The \texttt{caret} package can be used to provide a better confusion
matrix.

\begin{Shaded}
\begin{Highlighting}[]
\CommentTok{\# Round predictions to obtain binary values (0 or 1)}
\NormalTok{roundedresults }\OtherTok{\textless{}{-}} \FunctionTok{sapply}\NormalTok{(results, round, }\AttributeTok{digits=}\DecValTok{0}\NormalTok{)}
\NormalTok{roundedresultsdf }\OtherTok{\textless{}{-}} \FunctionTok{data.frame}\NormalTok{(roundedresults)}

\CommentTok{\# Convert variables to factors for proper processing in \textasciigrave{}caret\textasciigrave{}}
\NormalTok{roundedresultsdf}\SpecialCharTok{$}\NormalTok{actual }\OtherTok{\textless{}{-}} \FunctionTok{as.factor}\NormalTok{(roundedresultsdf}\SpecialCharTok{$}\NormalTok{actual)}
\NormalTok{roundedresultsdf}\SpecialCharTok{$}\NormalTok{prediction }\OtherTok{\textless{}{-}} \FunctionTok{as.factor}\NormalTok{(roundedresultsdf}\SpecialCharTok{$}\NormalTok{prediction)}

\CommentTok{\# Load caret package (install if not available)}
\ControlFlowTok{if}\NormalTok{ (}\SpecialCharTok{!}\FunctionTok{require}\NormalTok{(caret)) }\FunctionTok{install.packages}\NormalTok{(}\StringTok{"caret"}\NormalTok{, }\AttributeTok{dependencies=}\ConstantTok{TRUE}\NormalTok{)}
\FunctionTok{library}\NormalTok{(caret)}

\CommentTok{\# Generate confusion matrix with advanced metrics}
\NormalTok{conf\_matrix }\OtherTok{\textless{}{-}} \FunctionTok{confusionMatrix}\NormalTok{(roundedresultsdf}\SpecialCharTok{$}\NormalTok{prediction, roundedresultsdf}\SpecialCharTok{$}\NormalTok{actual)}

\CommentTok{\# Print the confusion matrix and performance metrics}
\FunctionTok{print}\NormalTok{(conf\_matrix)}
\end{Highlighting}
\end{Shaded}

\begin{verbatim}
Confusion Matrix and Statistics

          Reference
Prediction  0  1
         0 35  7
         1  1 24
                                         
               Accuracy : 0.8806         
                 95% CI : (0.7782, 0.947)
    No Information Rate : 0.5373         
    P-Value [Acc > NIR] : 1.956e-09      
                                         
                  Kappa : 0.7566         
                                         
 Mcnemar's Test P-Value : 0.0771         
                                         
            Sensitivity : 0.9722         
            Specificity : 0.7742         
         Pos Pred Value : 0.8333         
         Neg Pred Value : 0.9600         
             Prevalence : 0.5373         
         Detection Rate : 0.5224         
   Detection Prevalence : 0.6269         
      Balanced Accuracy : 0.8732         
                                         
       'Positive' Class : 0              
                                         
\end{verbatim}

\section{\texorpdfstring{Full workflow using
\texttt{Caret}}{Full workflow using Caret}}\label{full-workflow-using-caret}

The \texttt{Caret} package offers an alternative to \texttt{neuralnet}
to perform a more efficient process.

\begin{Shaded}
\begin{Highlighting}[]
\CommentTok{\# Load required libraries (install if necessary)}
\ControlFlowTok{if}\NormalTok{ (}\SpecialCharTok{!}\FunctionTok{require}\NormalTok{(caret)) }\FunctionTok{install.packages}\NormalTok{(}\StringTok{"caret"}\NormalTok{, }\AttributeTok{dependencies=}\ConstantTok{TRUE}\NormalTok{)}
\ControlFlowTok{if}\NormalTok{ (}\SpecialCharTok{!}\FunctionTok{require}\NormalTok{(nnet)) }\FunctionTok{install.packages}\NormalTok{(}\StringTok{"nnet"}\NormalTok{, }\AttributeTok{dependencies=}\ConstantTok{TRUE}\NormalTok{)  }\CommentTok{\# Needed for neural networks}
\FunctionTok{library}\NormalTok{(caret)}
\FunctionTok{library}\NormalTok{(nnet)}
\end{Highlighting}
\end{Shaded}

Data can be normalized using using caret's \texttt{preProcess} function

\begin{Shaded}
\begin{Highlighting}[]
\NormalTok{preproc }\OtherTok{\textless{}{-}} \FunctionTok{preProcess}\NormalTok{(mydata, }\AttributeTok{method =} \FunctionTok{c}\NormalTok{(}\StringTok{"range"}\NormalTok{))  }\CommentTok{\# Min{-}Max normalization}
\NormalTok{normData }\OtherTok{\textless{}{-}} \FunctionTok{predict}\NormalTok{(preproc, mydata)}
\end{Highlighting}
\end{Shaded}

Test/Train splitting is automated

\begin{Shaded}
\begin{Highlighting}[]
\CommentTok{\# Split data into training (67\%) and testing (33\%) sets}
\FunctionTok{set.seed}\NormalTok{(}\DecValTok{12345}\NormalTok{)}
\NormalTok{trainIndex }\OtherTok{\textless{}{-}} \FunctionTok{createDataPartition}\NormalTok{(normData}\SpecialCharTok{$}\NormalTok{dividend, }\AttributeTok{p =} \FloatTok{0.67}\NormalTok{, }\AttributeTok{list =} \ConstantTok{FALSE}\NormalTok{)}
\NormalTok{trainset }\OtherTok{\textless{}{-}}\NormalTok{ normData[trainIndex,]}
\NormalTok{testset }\OtherTok{\textless{}{-}}\NormalTok{ normData[}\SpecialCharTok{{-}}\NormalTok{trainIndex,]}
\end{Highlighting}
\end{Shaded}

Model variability can be managed by crossvalidation

\begin{Shaded}
\begin{Highlighting}[]
\NormalTok{ctrl }\OtherTok{\textless{}{-}} \FunctionTok{trainControl}\NormalTok{(}\AttributeTok{method =} \StringTok{"cv"}\NormalTok{, }\AttributeTok{number =} \DecValTok{10}\NormalTok{)  }\CommentTok{\# 10{-}fold cross{-}validation}
\end{Highlighting}
\end{Shaded}

\begin{Shaded}
\begin{Highlighting}[]
\CommentTok{\# Train a neural network model using caret\textquotesingle{}s \textasciigrave{}train\textasciigrave{} function}
\FunctionTok{set.seed}\NormalTok{(}\DecValTok{12345}\NormalTok{)}
\NormalTok{nn\_model }\OtherTok{\textless{}{-}} \FunctionTok{train}\NormalTok{(dividend }\SpecialCharTok{\textasciitilde{}}\NormalTok{ ., }\AttributeTok{data =}\NormalTok{ trainset,}
                  \AttributeTok{method =} \StringTok{"nnet"}\NormalTok{,  }\CommentTok{\# Use neural network model}
                  \AttributeTok{trControl =}\NormalTok{ ctrl,  }\CommentTok{\# Cross{-}validation}
                  \AttributeTok{tuneGrid =} \FunctionTok{expand.grid}\NormalTok{(}\AttributeTok{size =} \FunctionTok{c}\NormalTok{(}\DecValTok{4}\NormalTok{, }\DecValTok{2}\NormalTok{), }\AttributeTok{decay =} \FloatTok{0.01}\NormalTok{),  }\CommentTok{\# Hidden layers, weight decay}
                  \AttributeTok{maxit =} \DecValTok{1000}\NormalTok{,  }\CommentTok{\# Max iterations}
                  \AttributeTok{trace =} \ConstantTok{FALSE}\NormalTok{)  }\CommentTok{\# Suppress output during training}

\CommentTok{\# Print model summary}
\FunctionTok{print}\NormalTok{(nn\_model)}
\end{Highlighting}
\end{Shaded}

\begin{verbatim}
Neural Network 

134 samples
  5 predictor

No pre-processing
Resampling: Cross-Validated (10 fold) 
Summary of sample sizes: 120, 120, 121, 121, 120, 120, ... 
Resampling results across tuning parameters:

  size  RMSE       Rsquared   MAE      
  2     0.2009342  0.8309003  0.1048632
  4     0.2007985  0.8313619  0.1021993

Tuning parameter 'decay' was held constant at a value of 0.01
RMSE was used to select the optimal model using the smallest value.
The final values used for the model were size = 4 and decay = 0.01.
\end{verbatim}

Once the model is trained we can proceed as usually

\begin{Shaded}
\begin{Highlighting}[]
\CommentTok{\# Make predictions on test set}
\NormalTok{y\_pred }\OtherTok{\textless{}{-}} \FunctionTok{predict}\NormalTok{(nn\_model, testset)}

\CommentTok{\# Convert predictions to binary (0 or 1)}
\NormalTok{y\_pred\_bin }\OtherTok{\textless{}{-}} \FunctionTok{ifelse}\NormalTok{(y\_pred }\SpecialCharTok{\textgreater{}} \FloatTok{0.5}\NormalTok{, }\DecValTok{1}\NormalTok{, }\DecValTok{0}\NormalTok{)}

\CommentTok{\# Convert actual values and predictions to factors for confusion matrix}
\NormalTok{testset}\SpecialCharTok{$}\NormalTok{dividend }\OtherTok{\textless{}{-}} \FunctionTok{as.factor}\NormalTok{(testset}\SpecialCharTok{$}\NormalTok{dividend)}
\NormalTok{y\_pred\_bin }\OtherTok{\textless{}{-}} \FunctionTok{as.factor}\NormalTok{(y\_pred\_bin)}

\CommentTok{\# Generate confusion matrix and performance metrics}
\NormalTok{conf\_matrix }\OtherTok{\textless{}{-}} \FunctionTok{confusionMatrix}\NormalTok{(y\_pred\_bin, testset}\SpecialCharTok{$}\NormalTok{dividend)}

\CommentTok{\# Print confusion matrix and evaluation metrics}
\FunctionTok{print}\NormalTok{(conf\_matrix)}
\end{Highlighting}
\end{Shaded}

\begin{verbatim}
Confusion Matrix and Statistics

          Reference
Prediction  0  1
         0 34  4
         1  1 27
                                         
               Accuracy : 0.9242         
                 95% CI : (0.832, 0.9749)
    No Information Rate : 0.5303         
    P-Value [Acc > NIR] : 3.522e-12      
                                         
                  Kappa : 0.8471         
                                         
 Mcnemar's Test P-Value : 0.3711         
                                         
            Sensitivity : 0.9714         
            Specificity : 0.8710         
         Pos Pred Value : 0.8947         
         Neg Pred Value : 0.9643         
             Prevalence : 0.5303         
         Detection Rate : 0.5152         
   Detection Prevalence : 0.5758         
      Balanced Accuracy : 0.9212         
                                         
       'Positive' Class : 0              
                                         
\end{verbatim}

Notice that the model shows slightly better than the first one, probably
due to the fact that it has been minimally optimized.




\end{document}
